\IfFileExists{IEEEtran.cls}{%
  \documentclass[10pt,journal,compsoc]{IEEEtran}
}{%
  % The repository's minimal CI image may not contain IEEEtran.  This fallback
  % keeps the source syntax-checkable; camera-ready builds must use IEEEtran.
  \documentclass[10pt,twocolumn]{article}
}

\usepackage[T1]{fontenc}
\usepackage{amsmath,amssymb}
\usepackage{booktabs}
\usepackage{array}
\newcolumntype{L}[1]{>{\raggedright\arraybackslash}p{#1}}
\usepackage{cite}
% Expansion requires scalable Type-1 fonts that are absent from some minimal
% TeX images; protrusion alone keeps both the fallback and IEEE builds portable.
\usepackage[expansion=false]{microtype}
\usepackage{xcolor}
\usepackage{xspace}
\usepackage{url}
\usepackage[hidelinks]{hyperref}

\providecommand{\IEEEPARstart}[2]{#1#2}
\makeatletter
\@ifundefined{IEEEkeywords}{%
  \newenvironment{IEEEkeywords}{\par\small\noindent\textbf{Index Terms---}}{\par}
}{}
\makeatother

\newcommand{\system}{TaskGate\xspace}
\newcommand{\code}[1]{\texttt{\detokenize{#1}}}
\newcommand{\notmeasured}{\textcolor{gray}{not measured}}
\newcommand{\vect}[1]{\boldsymbol{#1}}
\newcommand{\R}{\mathbb{R}}
\newcommand{\N}{\mathbb{N}}
\newcommand{\grantlang}{\mathcal{L}}
\newcommand{\claimstatus}{Artifact draft; see Sec.~\ref{sec:evaluation}}

\title{TaskGate: Human-Approved, Budget-Constrained Authorization\\
for Autonomous Database Agents}

\author{Anonymous Author(s)%
\thanks{Anonymous manuscript prepared for review. The artifact status in this
draft is reported explicitly in Sec.~\ref{sec:evaluation}; unexecuted campaigns
are not presented as results.}}
\date{}

\begin{document}
\maketitle

\begin{abstract}
Autonomous agents often need a sequence of related database queries, creating a
choice between a broad reusable credential and repeated human approval.  We
present \system, a task-bound authorization and continuous usage-control design
for a deliberately narrow setting: restricted, read-only PostgreSQL~16
\code{SELECT} queries.  An authenticated subject proposes a versioned manifest
containing products, columns, mandatory row scopes, catalog and datasource
digests, expiry, and cumulative query, result-row, and database-time budgets.  An approval
service may reject or monotonically narrow that manifest.  A separate Go
gateway then mediates every query, validates the PostgreSQL abstract syntax
tree, injects scope predicates and an outer row limit, and reserves budget in a
control PostgreSQL database before contacting a credential-isolated business
database.  Durable request identifiers and conservative treatment of
indeterminate outcomes prevent automatic retries from consuming authority
twice.  Signed approval and query receipts bind decisions to protocol objects;
their signatures authenticate the issuing services, not a person's legal
intent.  We give a TLA+ model and an evaluation design covering policy bypass,
fault injection, and four performance baselines.  This artifact draft does not
report a completed full benchmark, an accepted 24-CPU-hour fuzz campaign, or a
human-subject study.
Its only approval-effort result is an algebraic count of approval events.  The
contribution is therefore a concrete, auditable enforcement design and
reproducibility contract, not a claim about measured performance or human
cognitive load.
\end{abstract}

\begin{IEEEkeywords}
agent authorization, database security, usage control, human approval,
PostgreSQL, idempotency, audit receipts
\end{IEEEkeywords}

\section{Introduction}
\label{sec:introduction}

\IEEEPARstart{A}{n} autonomous agent that can query a database is both a useful
planner and an untrusted query generator.  A long-lived database credential
authorizes much more than a user's current task.  Conversely, asking a human to
approve every query makes a multi-query investigation cumbersome and does not
solve retry, concurrency, or crash-accounting problems.  The missing mechanism
in our setting is not another authentication scheme: it is a task-bound grant
whose permitted query language and cumulative consumption remain enforced
after authentication and approval.

\system addresses this problem for one intentionally constrained deployment.
An external gateway is the only component given the business database
credential.  The agent requests logical data products, columns, mandatory
scope values, and a bounded lifetime.  A human-facing office-automation (OA)
service displays this structured request and may approve, reject, or narrow it.
After approval, the same grant governs every query in the task.  The gateway
does not ask an LLM whether a query ``matches the objective.''  It instead
checks deterministic properties of a PostgreSQL abstract syntax tree (AST),
rewrites the query through approved reporting views, and enforces a durable
vector budget.

The design separates three questions that are often conflated:

\begin{enumerate}
  \item \emph{Who is calling?}  OAuth/OIDC or an enterprise identity system is
  assumed to authenticate the human and agent workload.  The prototype uses
  static local identities only as a demo substitute.
  \item \emph{What future operations were approved?}  A versioned,
  RFC~8785-canonicalized manifest and a narrowing-only grant define the task
  envelope; the OA-signed approval receipt binds their digests.
  \item \emph{May this operation execute now?}  The gateway checks task state,
  catalog/schema agreement, SQL structure, remaining budget, and durable retry
  state immediately before execution.
\end{enumerate}

This division is important under prompt injection.  A compromised agent may
misread an objective or choose a malicious tool call, as illustrated by agent
security benchmarks such as AgentDojo~\cite{debenedetti2024agentdojo}.  \system
does not claim to repair the model's reasoning.  It limits the effects of that
reasoning to a human-visible, machine-enforceable database envelope.

The paper makes four scoped contributions:

\begin{itemize}
  \item a versioned authorization object and partial order under which OA
  decisions can only remove products, columns, scope values, time, sensitivity,
  or budget;
  \item an execution protocol that combines PostgreSQL AST enforcement,
  scope/limit rewriting, one in-flight query per task, reserve--execute--settle
  accounting, and durable request idempotency;
  \item approval and query evidence bound to manifest, grant, catalog, request,
  result, and audit-chain digests, with deliberately restrained integrity
  claims; and
  \item a formal safety model and a source-backed evaluation plan whose absent
  measurements remain visibly absent rather than being replaced with estimates.
\end{itemize}

We make no priority claim.  The design specializes ideas from task-based and
usage control, structured authorization, database query mediation, and recent
agent-authorization protocols.  Its value is the composition and the explicit
failure semantics for a concrete read-only database path.

\section{Related Work}
\label{sec:related}

\subsection{Structured and Task-Oriented Authorization}

OAuth Rich Authorization Requests (RAR) adds an
\code{authorization_details} parameter for fine-grained structured
authorization data~\cite{rfc9396}.  It conveys those details through the
authorization flow to a resource server, while leaving application-specific
detail types, comparison logic, and domain-specific processing to the
deployment.  \system supplies one such domain model together with a stateful
cumulative ledger, query rewriter, and crash/retry semantics; it does not
replace OAuth authentication or token issuance.

Task-Based Authorization Controls (TBAC) connect permissions and their usage to
the progress of enterprise tasks~\cite{thomas1998tbac}.  UCON$_{ABC}$ broadens
access control with authorizations, obligations, conditions, continuity, and
mutable attributes~\cite{park2004ucon}.  \system adopts their task lifecycle
and continuing-decision perspective but is much narrower: it has no general
obligation language, delegation administration, or arbitrary mutable subject
attributes.  Its mutable state is an explicit resource ledger and request
record for one read-only SQL task.

\subsection{Database Query Mediation}

Blockaid intercepts web-application queries and uses SMT reasoning to check
compliance with view-based data-access policies while preserving application
semantics~\cite{zhang2022blockaid}.  \system instead chooses a closed,
restricted SQL subset and deterministically rewrites approved logical products
to reporting views with mandatory predicates and a result cap.  This sacrifices
compatibility with arbitrary applications in exchange for a smaller policy
language, and adds human task grants, cumulative budgets, retry records, and
approval evidence.  Native views, role privileges, and read-only transactions
remain defense-in-depth rather than competitors to the gateway.

\subsection{Authorization and Guardrails for Agents}

Task Shield uses model-mediated task-alignment checks to identify instructions
and tool calls that do not serve a user's objective~\cite{jia2025taskshield}.
It addresses semantic alignment and indirect prompt injection.  \system is
complementary: it makes no semantic judgment, but continues to enforce a
structured database boundary even if the model-side shield fails.

PAuth derives natural-language slices that characterize task-implied calls and
uses signed envelopes to bind concrete operands to symbolic provenance
~\cite{sharma2026pauth}.  This promises operation-level precision for tasks
whose semantics can be correctly compiled.  \system declines that trust
assumption: natural-language objectives are displayed and logged but never
compiled into authority.  A human or policy route explicitly selects products,
columns, scopes, and budgets, after which enforcement is deterministic.

The Agentic Payment Protocol (AP2) uses trusted surfaces, signed mandates, and
linked checkout/payment evidence for agent-performed payments
~\cite{ap2v02}.  Its domain is commerce and payment, including payment-specific
roles and dispute evidence.  \system borrows the useful separation between an
agentic proposer and deterministic verification, but it neither handles payment
nor claims that its service signatures provide the legal non-repudiation that a
payment ecosystem may require.

The recent bounded capability receipts Internet-Draft combines a signed closed
scope with a durable reserve--execute--commit store, globally unique operation
identifiers, replay refusal, indeterminate outcomes, reconciliation, and
narrowing delegation~\cite{schrock2026bounded}.  It is an individual,
provisional, non-endorsed work in progress and the closest protocol-level
comparison.  \system is a non-conforming, domain-specific
analogue inspired by its reserve--execute--commit pattern, not an implementation
of the draft.  It replaces the draft's scalar budget with a vector-valued
database budget (queries, rows, and database milliseconds), omits holder proof
and child-capability lineage, and permits only one gateway and one in-flight
query per task.  Conversely, it specifies SQL-language enforcement,
catalog/schema attestation, OA presentation, and SQL-specific result receipts
beyond the draft's generic operation-evidence interface.

\begin{table*}[t]
\centering
\caption{Relationship to the most direct precedents. ``Not the focus'' avoids
inferring that a cited system could not be extended.}
\label{tab:related}
\footnotesize
\setlength{\tabcolsep}{4pt}
\begin{tabular}{L{0.14\textwidth}L{0.20\textwidth}L{0.25\textwidth}L{0.31\textwidth}}
\toprule
Work & Primary policy/evidence unit & Principal enforcement focus & Relationship to \system \\
\midrule
OAuth RAR~\cite{rfc9396} & Structured OAuth authorization detail & Transport and processing of fine-grained authorization data & Can carry a task detail; does not itself define the SQL language, cumulative ledger, or crash semantics. \\
TBAC~\cite{thomas1998tbac} / UCON~\cite{park2004ucon} & Task step / continuing usage decision & Lifecycle, usage, obligations, conditions, mutable attributes & Conceptual foundation; \system is a concrete, restricted database instantiation. \\
Blockaid~\cite{zhang2022blockaid} & View-based data policy & Per-query semantic policy compliance for web applications & Stronger application compatibility; task approval and durable aggregate consumption are not its focus. \\
Task Shield~\cite{jia2025taskshield} & User Task Set and LLM \emph{ContributesTo} decision & Model-mediated semantic alignment under prompt injection & Complementary semantic defense; \system provides deterministic effect containment without judging meaning. \\
PAuth~\cite{sharma2026pauth} & Signed natural-language task, server-derived slices, and signed operand envelopes & Task-implied server calls & Derives precise authority from language; \system requires explicit structured authority and supports a narrower SQL domain. \\
AP2~\cite{ap2v02} & Signed commerce/payment mandate & Agentic checkout, payment verification, and evidence & Similar trusted-surface separation, but domain and effect semantics differ. \\
Bounded capability receipts~\cite{schrock2026bounded} & Provisional signed capability plus shared spend state & Scoped reserve--execute--commit, replay refusal, and reconciliation & Non-endorsed individual draft; \system is a non-conforming database analogue with a vector budget and without holder proof or child lineage. \\
\bottomrule
\end{tabular}
\end{table*}

\section{Threat Model}
\label{sec:threat}

\subsection{System and Trust Boundary}

The system contains an agent client, an external Go gateway, an OA approval
service, a control PostgreSQL database, and a separate business PostgreSQL~16
database.  The gateway is the only agent-reachable component holding the
business reader credential.  The control database persists task state, grants,
budgets, request records, encrypted results, and audit events.  The business
database exposes only cataloged reporting views to a non-owner,
non-superuser, non-\code{BYPASSRLS} role.

We trust the gateway binary and its signing key, OA's authentication and signing
service, the two PostgreSQL engines, reporting-view definitions, the catalog,
and administrators who can modify those components.  We also trust the
enterprise identity provider assumed to precede \system.  The prototype's
static tokens and local passwords are demonstration mechanisms, not this
identity assumption's production realization.

\subsection{Adversary}

The agent may be fully compromised.  It may submit malformed or adversarial
SQL, multi-statements, nested queries, aliases, comments, dangerous functions,
unexpected operators, repeated calls, colliding request identifiers, and calls
after expiry or revocation.  It may exploit concurrency and may disconnect at
any point.  Business data returned by an earlier query may itself contain a
prompt injection.  A network attacker may replay an observed OA callback, but
does not know the callback HMAC secret or signing keys.  Components may crash
between reservation, connector invocation, result persistence, and settlement.

The adversary does not control the host, a database superuser, catalog
publisher, OA signing service, gateway signing service, or an approving human.
Compromise of any of these trusted components can authorize, execute, or rewrite
events outside the model.  Denial of service, traffic analysis, microarchitectural
side channels, and vulnerabilities in PostgreSQL or cryptographic libraries are
also outside the proved scope.

\subsection{Security Goals}

For the manual route, no query should reach the business connector without a
valid OA record of the authenticated human's approval.  For any executing query
$q$ and grant $G$, the gateway should ensure $q\in\grantlang(G)$, as defined in
Sec.~\ref{sec:model}.  Across each task and each resource dimension $j$,

\begin{equation}
  used_j + reserved_j \le limit_j.
  \label{eq:budget-invariant}
\end{equation}

The same $(task\_id,request\_id)$ and request digest must not trigger a second
connector invocation or second charge; binding the identifier to a different
request must fail closed.  Revocation, expiry, terminal state, and catalog/schema
drift must prevent new queries.  Business database credentials and uncataloged
base tables must remain unavailable to the agent.

These are authorization and accounting goals, not data non-interference.
\system does not bound inference from permitted results, provide differential
privacy, prove that the declared objective is truthful, or guarantee the
availability of a query that is authorized.

\paragraph{Manual and automatic routes.}
The paper's ``human-approved'' claim applies only when the catalog selects the
manual OA route.  The demonstration catalog also contains an explicitly labeled
low-sensitivity automatic route.  On that route, submitting the draft is an
authenticated requester action and OA issues its approval receipt with that
requester as the attributed actor.  This records requester assent and a policy
decision, but it does not establish review by a separate human approver.
Security experiments intended to substantiate such review must select the
manual route.

\section{Authorization Model}
\label{sec:model}

\subsection{Manifest and Grant}

Let a manifest $M$ contain

\begin{align*}
M = (&version, task, human, agent, objective,\\
     &P, C, S, sensitivity, \vect{b}, catalogVersion,\\
     &catalogHash, callbackContext, nonce),
\end{align*}

where $P$ is a nonempty set of logical products, $C(p)$ is a nonempty set of
columns for every $p\in P$, $S$ is a complete map of mandatory scope
constraints, and

\begin{equation}
\vect{b}=(b_q,b_r,b_t,b_{timeout},b_{ttl})
\end{equation}

contains maximum queries, result rows, cumulative database milliseconds,
per-query timeout, and task lifetime.  The gateway derives $human$ and $agent$
from the authenticated principal; neither is accepted as an agent-supplied
manifest field.

Sets are sorted before serialization.  The manifest digest is

\begin{equation}
\begin{aligned}
 H_M=\operatorname{SHA256}(&\texttt{TASKGATE-MANIFEST-V1\textbackslash0}\\
       &\parallel \operatorname{JCS}(M)),
\end{aligned}
 \label{eq:manifest-digest}
\end{equation}

where JCS is RFC~8785 canonical JSON~\cite{rfc8785}.  The domain separator
prevents this digest from being confused with a grant or receipt digest.  Fixed
cross-component vectors are part of the artifact contract.

The OA selects a grant core $G$ that repeats identity and provenance fields and
replaces requested products, columns, scope, sensitivity, expiry, and budget
with approved values.  A candidate $G'$ is a valid narrowing of $G$, written
$G'\preceq G$, only if:

\begin{align}
 P'&\subseteq P, & C'(p)&\subseteq C(p),\\
 S'&\preceq_S S, & \vect{b}'&\leq\vect{b},\\
 expiry'&\leq expiry, & sensitivity'&\leq sensitivity,
 \label{eq:narrowing}
\end{align}

and all identity, task, objective, catalog, and manifest-digest fields are
equal.  For an enum scope, $S'\preceq_S S$ means set inclusion; for a date
range, it means interval containment.  Unknown or removed mandatory scope keys
are rejected.  The component-wise order makes narrowing transitive and rules
out approval-time expansion.

An \code{ApprovalReceiptV1} contains a receipt and task identifier, decision,
$H_M$, approved grant-core digest, approver identifier, issue time, key
identifier, and Ed25519 signature.  Ed25519 is specified by RFC~8032
~\cite{rfc8032}.  The callback's HMAC authenticates transport and supports
replay control; the Ed25519 receipt persistently authenticates the OA service's
record.  A final \code{TaskGrantV1} combines the core and receipt only after the
gateway verifies the receipt and the narrowing relation.

\subsection{Authorized Query Language}

Let $A(q)$ be the PostgreSQL parse tree of agent SQL $q$.  For a valid grant
$G$, $q\in\grantlang(G)$ only if all of the following hold:

\begin{enumerate}
  \item $A(q)$ contains exactly one top-level \code{SELECT}; write statements,
  utility statements, locking forms, and disallowed constructs are absent;
  \item each \code{SELECT}, CTE, and subquery has its own lexical relation
  scope.  A relation name denotes either an unqualified logical product in
  $P_G$, a CTE, or a derived relation.  Aliases bind to that concrete product
  or derived relation.  A qualified column must exist in the named relation;
  an unqualified column is accepted only when it resolves to exactly one
  in-scope relation.  Ambiguous names, stars, and constructs that can expand
  the approved projection are rejected;
  \item the implementation materializes a product-indexed authorization
  environment
  $E_G(p)=(R(p), C(p), F(p), A(p), O(p))$ for relation, approved columns,
  scalar functions, aggregates, and operators.  Derived columns carry the set
  of source products used to compute them;
  \item a function, aggregate, or catalog-controlled operator over one product
  must be in that product's allowlist.  An expression over multiple products
  must use names in the intersection of the source products' allowlists.
  Constant expressions that reference no data columns use a separate
  task-global safe list.  Structural SQL nodes such as \code{AND}, \code{OR},
  \code{NULL TEST}, and \code{CASE} are syntax, not catalog-controlled
  functions or operators;
  \item system catalogs, schema-qualified agent objects, and dangerous
  introspection or side-effecting functions are absent; and
  \item deterministic rewriting replaces each logical product with a CTE over
  its cataloged reporting view, exposing only approved columns and applying all
  $S_G$ predicates, then wraps the outer query with a limit no greater than the
  remaining row budget.
\end{enumerate}

The definition is syntactic and catalog-relative.  It says neither that $q$
serves $objective$ nor that the reporting view itself implements the
organization's intended policy.  Catalog and schema correctness remain trusted
inputs checked for drift, not inferred policy.

\subsection{Budget and Request Semantics}

For the consumable resource vector, use
$\vect{l}=(l_q,l_r,l_t)$, $\vect{u}=(u_q,u_r,u_t)$, and
$\vect{r}=(r_q,r_r,r_t)$.  A reservation is allowed atomically only if the task
is active and

\begin{equation}
  \vect{u}+\vect{r}+\vect{a}\leq\vect{l},
  \qquad \vect{a}=(1,a_r,a_t).
  \label{eq:reserve}
\end{equation}

The control store locks the task and budget rows and requires $r_q=0$, which is
the implementation of one in-flight query per task.  On a known pre-connector
failure it releases $\vect{a}$.  On success it moves bounded actual use
$(1,x_r,x_t)\leq\vect{a}$ from reserved to used.  If the connector may have
executed but the outcome cannot be established, it charges the full reservation
and records \code{INDETERMINATE}.  Conservative charging preserves
Eq.~\eqref{eq:budget-invariant} but may reduce availability.

The unique key $(task\_id,request\_id)$ is bound to a digest of the tool name,
task, query or plan, and output options.  Reuse with the same digest returns the
first durable result or terminal status.  Reuse with a different digest is a
conflict.  An agent can intentionally choose a new request identifier for the
same SQL; that is a new budgeted operation, not a violation of retry
idempotency.

Revocation makes the task terminal for new reservations.  It does not cancel a
query that has already been reserved; that query remains bounded by its already
approved timeout and grant expiry.

\section{Design}
\label{sec:design}

\subsection{Architecture and Isolation}

Figure~\ref{fig:architecture} shows the fixed deployment route.  Control and
business state are placed in different PostgreSQL instances, roles, credentials,
and volumes.  The agent reaches the gateway through MCP tools but has no network
route or credential for Business PG.  OA receives authorization drafts and
returns signed decisions.  Control PG is the serialization point for task state,
request idempotency, budgets, and audit order.

\begin{figure*}[t]
\centering
\fbox{\begin{minipage}{0.94\textwidth}
\centering
Untrusted Agent / MCP Client
$\longrightarrow$
\textbf{Go Gateway}
$\longrightarrow$
Business PostgreSQL~16 (reporting views, restricted reader)\\[2mm]
\hfill
$\searrow$ OA approval service (display, approve/reject/narrow, sign)
\qquad
$\searrow$ Control PostgreSQL (task, grant, budget, request, result, audit)
\end{minipage}}
\caption{Fixed \system deployment.  The prototype and model assume one gateway;
different tasks may execute concurrently, but each task has at most one reserved
query.}
\label{fig:architecture}
\end{figure*}

The deployment's business role is \code{LOGIN NOSUPERUSER NOCREATEDB
NOCREATEROLE NOINHERIT NOREPLICATION NOBYPASSRLS}.  It receives \code{SELECT}
only on named reporting views and has read-only transactions and a server-side
statement timeout.  PostgreSQL documents that a read-only transaction rejects
ordinary data modification and schema commands~\cite{postgres16readonly}; this
is defense-in-depth because the AST policy rejects non-\code{SELECT} statements
first.  The default Compose topology does not publish the business database's
port to the host.  Debug exposure, if added, is a non-paper override.

\subsection{Approval Path}

The approval path has seven stages:

\begin{enumerate}
  \item The authenticated client lists cataloged products and explicitly
  supplies an objective, products, nonempty per-product columns, all mandatory
  scopes, and optional smaller budgets.
  \item The gateway resolves a catalog route and hard ceilings, derives identity
  fields, normalizes sets, creates $M$, and computes $H_M$.
  \item OA displays the objective, agent identity, products, columns, scope,
  sensitivity, every budget dimension, and lifetime.  The display is generated
  from the same manifest that is hashed, not from a separate summary.
  \item On the manual route, the approver rejects, accepts, or constructs a
  candidate grant core.  Narrowing can remove products, columns, and enum
  values; contract date ranges; shorten lifetime; or lower any budget.  The
  current UI keeps sensitivity fixed; the protocol validator permits a lower
  candidate ceiling but rejects an increase.  On the automatic route, OA uses
  the unchanged requested core when the authenticated requester submits it.
  Neither path can add authority.
  \item OA signs a receipt binding the decision to $H_M$ and the approved core
  digest.  The callback also carries an HMAC, event identifier, callback context,
  catalog version, and occurrence time.
  \item The gateway verifies HMAC freshness/replay, actor and task state, catalog
  and callback context, both digests, Ed25519 key identifier/signature, and
  $G'\preceq G_M$.
  \item One Control PG transaction records the approval event, immutable final
  grant, initial budget ledger, and new task state.
\end{enumerate}

Event-id replay returns the transactionally stored callback response.  An event
identifier reused for different bytes is rejected.  HMAC and Ed25519 have
separate jobs: compromise of a callback transport secret should not silently
change already retained receipt bytes, while a receipt alone is not accepted as
a fresh state transition.

\subsection{Execution Path}

Algorithm~\ref{alg:execution} gives the logical ordering.  The order is a
security property: in particular, the gateway completes policy analysis before
reserving, persists the reservation before contacting Business PG, and never
releases a reservation once it cannot establish that execution did not occur.

\begin{figure}[t]
\small
\begin{enumerate}
  \item Authenticate; load owned task and any existing request record.
  \item If the same request is durable, return it; if its digest differs, deny.
  \item Check active state, grant/TTL, catalog digest, and live datasource/schema readiness.
  \item Parse one PostgreSQL statement and validate the AST allowlists.
  \item Rewrite products, projections, scopes, and the outer row limit.
  \item Atomically reserve query, row, and database-time budget.
  \item Execute in a bounded read-only Business PG transaction.
  \item Atomically settle actual use and result, release known pre-execution
  failures, record post-execution \code{FAILED} outcomes, or conservatively mark
  ambiguous outcomes \code{INDETERMINATE}; append durable terminal audit
  evidence and persist a gateway-signed receipt before the same Control PG
  transaction commits.
\end{enumerate}
\caption{Security-relevant execution order.}
\label{alg:execution}
\end{figure}

Both \code{query\_sql} and the declarative \code{execute\_plan} require a
\code{request\_id}.  A QueryPlan compiler accepts product, projection,
aggregates, typed filters, ordering, and limit without SQL fragments, then sends
its output through the same AST policy as direct SQL.  Direct SQL can name only
logical products.  The rewriter creates internal CTE names unavailable to agent
SQL, selects approved physical view columns, adds mandatory predicates as
gateway-rendered literals, and applies a final cap.  No client string is
concatenated as a scope predicate or physical identifier.

The remaining grant lifetime bounds both statement timeout and the gateway
context.  Result collection is also capped in the connector.  Successful
results, post-execution failures, releases, indeterminate settlements, terminal
audit events, and signed receipts are persisted in the same Control PG
settlement transaction; transient settlement failure makes readiness fail and
enters an idempotent retry loop.
Startup recovery changes a left-over \code{RESERVED} record to
\code{INDETERMINATE} and charges its full reservation because the process cannot
reconstruct whether the connector crossed the read boundary.  Production startup
uses the configured gateway receipt key to persist that recovery receipt in the
same recovery transaction.

\subsection{Receipts and Audit}

A query receipt binds task, query and request identifiers; manifest, grant,
catalog, datasource, and schema digests; request digest and canonical SQL
fingerprint; policy decision; budget before, reservation, charge, and budget
after; row count and database time; result hash; terminal status/error;
operation and signing timestamps; audit sequence and previous audit hash; and a
gateway key identifier.  The gateway signs canonical, domain-separated V3
receipt bytes with Ed25519 and stores the signed receipt bytes in Control PG.
Repeated reads
return the stored receipt rather than re-signing a mutable query row.  The
receipt contains result metadata and a digest, not raw rows.  Receipt listing
uses a durable query-record cursor rather than a fixed prefix.  Audit-receipt
reads fetch the exact query audit events and return the terminal event,
predecessor evidence when present, the successor path from the terminal event
to the current audit checkpoint, and the checkpoint itself.
Archived tasks no longer accept new queries, but retained result ciphertext
remains readable to the owning principal until retention cleanup deletes the
encrypted result row or an administrator erases the result encryption key ID.
Ciphertext cleanup can be driven by a configured TTL scheduler or by an
administrator cutoff request, and active legal holds exclude a task's result
ciphertext from deletion.  Key-ID erasure marks the corresponding
\code{result\_encryption\_keys} row \code{ERASED}, appends audit evidence, and
causes later result reads to fail closed while leaving ciphertext rows, query
records, signed receipts, and audit evidence intact.  Disabling a principal
prevents tool listing and tool execution even if a client still presents the old
static bearer token.
The verifier keeps compatibility for older V1/V2 receipt formats while V3
binds \code{signed\_at} so historical gateway keys can have valid-from and
retirement windows.  The gateway publishes a
\code{taskgate-query-receipt-keyring/v1} bundle containing the active key ID,
historical verification keys, and validity windows, so external verifiers can
reconstruct the receipt verifier without private key material.  A shared
audit-chain verifier reconstructs the hash material for each returned event,
checks the predecessor and successor links, and binds the terminal event's
sequence and hashes to the signed receipt.  When configured with an external
anchor URL, the gateway also signs the current audit checkpoint as a
\code{taskgate-audit-checkpoint-anchor/v1} payload and posts it to an
independent log or WORM service.

For audit event fields $e_i=(id,task,query,actor,type,payload,time)$, the
implementation forms

\begin{equation}
 h_i=\operatorname{SHA256}(\operatorname{JSON}(h_{i-1},e_i)),
\end{equation}

where \code{JSON} is the fixed Go struct encoding with a normalized JSON
payload.  Unlike signed protocol objects, this audit encoding is not claimed to
be RFC~8785 or a cross-language wire format.

Control PG locks a single chain-head row while allocating the next sequence,
inserting the event, and updating the head.  This prevents benign concurrent
writers from creating forks.  Given a retained trusted head or signed receipt,
verification can reveal changed, removed, or reordered events on the observed
prefix.  It is not WORM storage: a privileged administrator who can replace the
database and recompute the entire unanchored chain can present an alternative
history.  Likewise, a query-receipt signature plus inclusion path proves that
the configured gateway key signed evidence for an event on this retained chain
head, not that Business PG independently attested execution.

\subsection{Catalog and Schema Drift}

The signed authorization includes catalog version, catalog SHA-256 digest,
datasource identifier, and reporting-schema digest.  Every catalog product also
gives the expected reporting-view column names, order, and PostgreSQL types. At
startup, readiness checks, and query execution before budget reservation, the
connector compares those expectations with the live datasource marker,
\code{current\_database()}, \code{current\_user}, PostgreSQL major version,
\code{information\_schema.columns}, and normalized \code{pg\_get\_viewdef}
output.  Any missing, extra, reordered, or mistyped column, or any changed
normalized reporting-view definition, makes readiness fail and blocks the
execution path.  Catalog upgrades therefore require draining or replacing
active grants; an old grant is not silently reinterpreted under a new schema.
Catalog V1 compares generic \code{information\_schema.data\_type} values and
rejects type strings carrying modifiers.  It therefore does not attest numeric
precision/scale, character length, domain identity, or array element type.

\section{Formal Analysis}
\label{sec:formal}

\subsection{State Machine}

The TLA+ suite~\cite{lamport2002specifying} contains six finite abstractions.
\code{TaskGate.tla} models one task with finite sets of request identifiers,
receipt identifiers, relations, columns, and scopes, together with a scalar
budget for the core lifecycle.  \code{VectorBudget.tla} is a separate split
model of the query, row, and database-millisecond budget ledger, including
hard-limit archival under weak fairness.  \code{SQLAuthorization.tla} models a
finite two-product SQL authorization abstraction with product--column
provenance, qualified and unqualified binding, constant-only safe lists, and
multi-product function/operator intersections.  \code{MultiTaskAudit.tla}
models finite multi-task interleavings, global audit sequence order, and
revocation/expiry races with in-flight requests.  \code{ReceiptAudit.tla}
models persisted terminal receipt field semantics and binding to the terminal
audit event.  \code{RecoveryLiveness.tla} models conservative recovery
convergence under weak fairness.  The core model's task state is

\begin{equation*}
\begin{aligned}
\{&\mathtt{PENDING},\mathtt{ACTIVE},\mathtt{REVOKED},\\
  &\mathtt{EXPIRED},\mathtt{COMPLETED},\mathtt{FAILED}\}.
\end{aligned}
\end{equation*}

Each request moves through \code{NEW}, \code{RESERVED}, \code{EXECUTING},
\code{RESULT\_KNOWN}, \code{SETTLING}, \code{RECOVERING}, and one of the
terminal states \code{SETTLED}, \code{RELEASED}, \code{FAILED}, or
\code{INDETERMINATE}.
In the core model, state variables record approval and catalog-validity flags,
a monotonically narrowed grant history, processed receipt identifiers and their
effects, scalar used/reserved budget, per-request grant and validity snapshots,
request state, and an abstract execution count.  Actions cover grant narrowing,
a fresh approval, approval replay, invalid approval rejection, reservation,
duplicate request attempts, connector/result/settlement phases, a definite
pre-connector release, three crash outcomes, revocation, expiry, rejection,
completion, and catalog drift.  The split models keep SQL authorization,
vector accounting, audit ordering, receipt semantics, and recovery liveness in
separate finite state spaces.

The suite deliberately does not combine all concerns into one state space.  It
does not model full PostgreSQL parser syntax, arbitrary AST nesting,
cryptographic signature bytes, JSON canonicalization bytes, row contents,
network I/O, external WORM/timestamp storage, or wall-clock scheduling.
Instead, grant containment, approval/catalog checks, parser output, audit hash
material, and recovery scheduling are abstracted at the boundary named by each
split model.  The core model permits an already-reserved request to finish
after a later terminal transition; its terminal invariant forbids \emph{new}
requests.  The repository includes \code{formal/REFINEMENT.md}, an audit map
from each modeled action to the current Go method, Control PostgreSQL
transaction or invariant, and test or fault-injection point.  That map is
traceability evidence, not a mechanized refinement proof.

\subsection{Safety Properties}

The submitted finite configurations check the following named invariants:

\begin{itemize}
  \item \code{TypeOK}: every variable remains in its declared finite domain.
  \item \code{NoExecutionWithoutApproval}: an abstract execution implies that the
  approval-valid flag was true when its request was reserved.
  \item \code{ExecutedWithinGrant}: every executed request's relation, column, and
  scope belong to the grant snapshot taken at reservation.
  \item \code{GrantMonotonic}: each grant-history element narrows the preceding
  relation, column, scope, and scalar-limit components.
  \item \code{BudgetSafety}: $used+reserved\leq limit$ for the scalar instance.
  \item \code{VectorBudgetSafety}: per-dimension $used+reserved\leq limit$ in
  the split vector-budget instance.
  \item \code{TerminalChargeBounded}, \code{ReleaseChargesNothing}, and
  \code{IndeterminateChargesFullReservation}: vector terminal transitions
  charge no more than the reservation, release without charge, and charge the
  full reservation for indeterminate outcomes.
  \item \code{NoReservationAfterArchive} and
  \code{EventualArchiveAtHardLimit}: after a vector hard limit is reached and
  reservations drain, weak fairness eventually archives the task and no new
  reservation is admitted.
  \item \code{AcceptedQueriesAuthorized},
  \code{AcceptedSourcesMatchColumns}, \code{ConstantsUseGlobalSafeList},
  \code{UnqualifiedColumnsAreUnique}, and
  \code{MultiProductUsesIntersection}: accepted abstract SQL shapes preserve
  product--column provenance, reject ambiguous unqualified columns, and require
  functions/operators to be in the relevant product intersection.
  \item \code{SingleInFlightPerTask}, \code{AuditIsLinear},
  \code{NoNewRequestAfterTaskTerminal}, and
  \code{TerminalRequestsHaveTerminalAudit}: finite multi-task interleavings
  preserve per-task one-in-flight admission, global audit order, terminal-task
  admission closure, and terminal audit evidence.
  \item \code{ReceiptExistsForTerminal}, \code{ReceiptBindsTerminalAudit},
  \code{ChargeWithinReservation}, \code{BudgetTransitionValid}, and
  \code{StatusFieldsValid}: persisted terminal receipts have legal
  budget/status/result/error fields and bind the terminal audit event.
  \item \code{RecoveryConverges}: under weak fairness for the recovery action,
  every finite recovering request eventually becomes terminal.
  \item \code{ApprovalReplaySafe}: each receipt identifier has at most one approval
  effect.
  \item \code{AtMostOnceExecution}: each request identifier executes at most once.
  \item \code{NoNewQueryAfterTerminal}: revoked, expired, completed, or failed tasks
  start no new request.
  \item \code{CatalogFailClosed}: every started request records that the abstract
  catalog-match flag was true at reservation.
  \item \code{SingleInFlight}: at most one request is in a nonterminal processing
  state.
\end{itemize}

Here \code{CatalogFailClosed} is intentionally narrower than its name might
suggest: it says that the abstract catalog-match flag was true at reservation;
the Go/schema refinement is not part of the model.  Likewise, the scalar model
adds and settles one unit per request.  The split vector, SQL, audit, receipt,
and recovery models check their own finite boundaries without proving that all
Go and PostgreSQL code refines a single composed TLA+ specification.

Three local inductive arguments explain the abstraction.  First, reservation
is the only action that raises \code{reserved}; its guard establishes
$used+reserved\leq limit$, and every modeled settle/release transition preserves
the inequality.  Second, narrowing is component-wise subset order for relation,
column, and scope sets plus non-increasing scalar limit.  Third, a request is
never returned to \code{NEW}, and the only execution-producing actions advance
it, so duplicate attempts cannot execute it again.  Connecting those arguments
to production behavior depends on Control PG atomicity and on all writers using
the modeled transactions.

TLC is an explicit-state checker over a finite configuration, not a proof for
unbounded tasks, budgets, or SQL syntax.  We therefore report the exact model,
configuration, tool version, explored-state counts, and invariant outcome only
when generated logs are present.  Table~\ref{tab:formal-status} is generated
from artifact status; this prose asserts no unobserved TLC result.


\begin{table}[t]
\centering
\caption{Generated finite-model status.  A missing raw log is not a pass.}
\label{tab:formal-status}
\footnotesize
\begin{tabular}{ll}
\toprule
Field & Generated value \\
\midrule
TLC outcome & passed \\
Tool & TLC 1.7.1 \\
Model & \parbox[t]{0.48\columnwidth}{formal/TaskGate.tla} \\
Config & \parbox[t]{0.48\columnwidth}{formal/TaskGate.cfg} \\
Checked at & 2026-07-21T17:38:38Z \\
States generated & 14824257 \\
Distinct states & 3255552 \\
Search depth & 18 \\
Evidence & \parbox[t]{0.48\columnwidth}{formal/results/tlc.log; TLC result, no-error completion marker, final statistics, and referenced inputs SHA-256 verified} \\
\bottomrule
\end{tabular}
\end{table}


\section{Implementation}
\label{sec:implementation}

The prototype consists of an external Go gateway, a Go OA demonstration
service, Control PostgreSQL, and Business PostgreSQL~16.  It exposes MCP tools
to list products, request and inspect tasks, execute a QueryPlan or direct SQL,
read results and budgets, list receipts, complete a task, and revoke it.  An
auditor principal can read receipt/audit metadata but not raw query rows.

The SQL policy uses \code{pg\_query\_go/v6}, a Go binding around code extracted
from PostgreSQL's parser~\cite{libpgquery}, to parse and deparse the server
dialect.  The analyzer rejects multiple statements; non-\code{SELECT} roots;
schema-qualified and system objects; row locks; unsafe CTE forms; projection
expansion; and non-allowlisted functions, aggregates, and operators.  Parser
acceptance alone never authorizes a query.  The connector starts an explicit
read-only transaction, sets a local statement timeout and search path, fetches
at most the reserved rows, commits after a successful read, and rolls back on
any earlier error.

Control PG uses row locks for task/budget serialization and a unique constraint
on $(task\_id,request\_id)$.  It stores budget snapshots with each request,
encrypted result bytes using AES-256-GCM under a recorded result key ID, a
plaintext result SHA-256 digest, and append-only grant/audit/receipt rows
protected from ordinary updates and deletes by triggers.  Result key registry
rows can transition from \code{ACTIVE} to \code{ERASED} but cannot be
reactivated or deleted through ordinary SQL.  Terminal query rows reject
deletes, authorization-evidence rewrites, and updates after settlement.  A
single audit-head row provides a total order.  These database
mechanisms protect against application bugs and concurrent requests, not a
control-database superuser.

Protocol objects are RFC~8785-canonicalized and domain-separated before hashing
or signing.  OA receipts are Ed25519-signed and verified by key identifier; the
gateway can load a keyring of OA verification keys with valid-from and
retirement times, so grants issued before a key's retirement can continue to
verify while later signatures under that key are rejected.
For any terminal query without a stored receipt, the gateway constructs, signs,
and then persists a receipt binding task, query and request identities;
manifest, grant, request, catalog, datasource, and schema digests; policy
decision; before/reserved/charged/after budget evidence; outcome and result
hash; signing time; and the corresponding audit-chain position.  Its dedicated
key is
loaded through \code{GATEWAY\_RECEIPT\_KEY\_ID} and
\code{GATEWAY\_RECEIPT\_PRIVATE\_KEY}; a still-reserved query exposes only an
unsigned status record.  HMAC remains on OA callbacks for freshness and
byte-level transport authentication.  Demo key derivation and environment-
variable secret storage are development conveniences, not a production
key-management design.

The implementation checks reporting schemas both on connector creation and on
readiness.  Compose runs the gateway and OA containers as non-root processes
with read-only root filesystems and \code{no-new-privileges}; only the gateway
network reaches Business PG by default.  These controls establish the fixed
paper deployment, not a general PostgreSQL wire proxy, extension, or
multi-database authorization platform.

\section{Evaluation}
\label{sec:evaluation}

\subsection{Questions and Reporting Discipline}

The evaluation is organized around four questions:

\begin{description}
  \item[RQ1:] Does the implementation reject queries, callbacks, receipts, and
  direct database accesses outside the stated model?
  \item[RQ2:] Do concurrency, retries, revocation, expiry, and injected crashes
  preserve authorization, budget, and at-most-once invariants?
  \item[RQ3:] What latency, throughput, CPU, memory, Control PG, and receipt
  storage overhead does each enforcement layer add?
  \item[RQ4:] How many approval events are required under a per-query policy
  versus a task-envelope policy for the same partition into tasks?
\end{description}

The full security, fuzz, formal, and performance campaigns are not silently
replaced by unit-test counts.  Generated tables ingest machine-readable raw
artifacts; without those files, they print \notmeasured.  No cell in this draft
is an estimate, illustrative benchmark, or manually transcribed result.

\subsection{Functional and Adversarial Method}

The security corpus covers multi-statements, comments, malformed parser edges,
CTEs, recursive and data-modifying forms, subqueries, aliases, set operations,
row locks, system catalogs, schema qualification, projection expansion,
disallowed functions/operators, and attempts to weaken or bypass mandatory
scope.  Protocol tests mutate selected authorization-relevant manifest, grant,
approval-receipt, and query-receipt fields and grant-expansion dimensions;
exercise callback replay and invalid signature or key identifiers; reuse a
request identifier with equal and unequal request digests; and verify that a
reported live-schema mismatch fails closed before reservation.

Fault injection is placed (i) before and after reservation, (ii) immediately
before connector invocation, (iii) after the query can have executed, and
(iv) before and after settlement/result persistence.  Expected outcomes are
release only when non-execution is known, bounded actual settlement when the
result is known, and full conservative charge with \code{INDETERMINATE}
otherwise.  Negative deployment tests attempt direct Business PG connection,
base-table access, and use of wrong credentials.  Separately, a deterministic
prompt-boundary corpus maps untrusted-text cases to representative SQL attempts
and checks policy rejection; it neither invokes a model nor establishes
prompt-injection robustness.

The long Docker fuzz campaign has three targets: arbitrary SQL-policy input must
not panic; formatting changes must preserve the policy outcome and canonical
allowed-query decision; and arbitrary QueryPlan compilation must not panic.  A
machine-readable campaign manifest names those targets and records SHA-256
digests for each combined per-target Go/GNU-time log.  The configured
\code{FUZZ\_CPU\_HOURS=24} is only a schedule request: acceptance requires the
parsed sum of actual user and system CPU seconds to reach 24 hours.  Fixed unit
and property vectors, rather than the long campaign, cover JCS canonicalization
and grant narrowing.  Docker integration tests exercise approved structured and
direct SQL plus database-role restrictions against PostgreSQL~16.  A scoped
Compose-backed differential test authorizes and deparses three representative
agent queries---an ordered projection, a CTE with a date literal, and a derived
table alias---executes each rewrite on PostgreSQL~16, and compares its ordered
rows with an independently written physical-view query.

The planned full security acceptance bar combines completion of those three
fuzz targets with no panic, stable external error codes, zero unauthorized
connector crossings in the defined corpus, and no budget-invariant violation.
Until machine-readable results and the checksummed raw logs record those facts,
the paper does not claim that the bar was met.

A generated \code{partial} outcome denotes harness-validation evidence only:
it does not establish the 24-CPU-hour threshold, unauthorized-connector-crossing
count, budget-fault invariant count, or full security acceptance.  Measured
cells and unmeasured cells remain separate in the generated table.


\begin{table}[t]
\centering
\caption{Generated security-campaign status.  A partial harness-validation run is not an accepted full campaign.}
\label{tab:security-status}
\footnotesize
\begin{tabular}{lr}
\toprule
Measure & Generated value \\
\midrule
Campaign outcome & partial \\
SQL corpus passed/cases & 15/15 \\
Prompt boundaries passed/cases & 2/2 \\
Unauthorized connector crossings & \notmeasured \\
Budget invariant violations & \notmeasured \\
Panics & 0 \\
Aggregate fuzz CPU-hours & 0.00185 \\
\midrule
\multicolumn{2}{>{\raggedright\arraybackslash}p{0.86\columnwidth}}{Evidence: partial; 127 checksummed security inputs; deterministically reconstructed from evaluation/security/results.json} \\
\bottomrule
\end{tabular}
\end{table}


\subsection{Performance Method}

The four baselines are (B1) direct PostgreSQL, (B2) native reporting views,
(B3) an AST-only gateway without approval, control ledger, encrypted result
persistence, or signed receipt, and (B4) full \system.  The current workload is
TPC-derived rather than a standards-compliant TPC-H/TPC-DS result.  Each
configuration records a baseline-order seed and executes cells in a seeded
random order at concurrency $1,8,32$ after warmup; the checked configs use
distinct-task concurrency so the experiment does not violate per-task
serialization.  Each cell requires at least 30 measured repetitions.

The driver records per-query throughput, p50/p95 end-to-end latency, p99 only
for rows with at least 10,000 observations, deterministic p50 bootstrap
intervals, direct-baseline latency ratios, gateway and database CPU, peak
resident memory, Control PG transaction counts/latency, receipt bytes, and
available baseline-specific component times.  For the full
gateway, current probes group authorization/control before connector invocation,
database execution, connector overhead, and post-connector settlement/receipt;
they do not yet isolate parse from reservation or settlement from encryption and
signing.  Environment artifacts include CPU model and allocation, memory,
storage, kernel, Docker, Go, PostgreSQL, image digests, database parameters,
dataset-generator version/seed and checksums, commit, warmup/cache rule, seed,
task-concurrency mode, and exact command.  Baselines use semantically
equivalent query variants, the same data, client host, connection policy, and
result consumption.  Raw JSON or CSV is the sole source for reported tables and
figures.


\begin{table*}[t]
\centering
\caption{Generated performance status.  No complete source-backed record set is present; no latency or throughput conclusion is drawn.}
\label{tab:performance-status}
\footnotesize
\begin{tabular}{lllllll}
\toprule
Experiment & SF & Concurrency & Baseline & Runs & Throughput & p50/p95/p99 \\
\midrule
\notmeasured & \notmeasured & \notmeasured & \notmeasured & \notmeasured & \notmeasured & \notmeasured \\
\bottomrule
\end{tabular}
\\[1mm]\scriptsize Generated reason: source-backed summary loaded.
\end{table*}


\subsection{Computable Approval Count}

Suppose an accepted workload is partitioned into $T$ approved tasks and task
$t$ contains $n_t\geq 1$ queries.  A policy that asks for approval before every
query creates

\begin{equation}
 A_{query}=\sum_{t=1}^{T}n_t
\end{equation}

approval events.  A manual \system task envelope, assuming one decision and no
revision loop per accepted task, creates

\begin{equation}
 A_{task}=T, \qquad
 A_{query}-A_{task}=\sum_{t=1}^{T}(n_t-1).
 \label{eq:approval-count}
\end{equation}

This is a count implied by the workflow, not a measured human outcome.  A
rejection, re-application, or organization that requires per-query step-up
approval adds events and must be counted from its trace.  The automatic demo
route includes an authenticated requester submission but no separate
approval-review event; it is excluded rather than treated as a zero-interaction
case.  We do not infer reduced fatigue, cognitive load, review time, or error
rate from Eq.~\eqref{eq:approval-count}.


\begin{table}[t]
\centering
\caption{Analytical approval-event count; these are formulas, not participant measurements.}
\label{tab:approval-count}
\footnotesize
\begin{tabular}{p{0.20\columnwidth}p{0.26\columnwidth}p{0.38\columnwidth}}
\toprule
Policy & Human approvals & Assumption \\
\midrule
Per query & $\sum_{t=1}^{T} n_t$ & one decision per query \\
TaskGate manual & $T$ & one accepted decision per task \\
Difference & $\sum_{t=1}^{T}(n_t-1)$ & same task partition \\
\bottomrule
\end{tabular}
\end{table}


\subsection{Reproducibility}
\label{sec:reproducibility}

The artifact contract exposes five repository commands:

\begin{itemize}\raggedright
  \item \code{make verify}: functional, race, security, and Compose integration
  tests;
  \item \code{make formal}: TLC with the submitted finite configurations and
  retained logs;
  \item \code{make eval-smoke}: a small four-baseline pipeline check;
  \item \code{make eval-full}: the complete security/performance campaigns; and
  \item \code{make paper}: regenerate tables/figures from raw artifacts and
  compile this source.
\end{itemize}

The public paper entrypoints first reconstruct the generated evaluation summary
from checksummed raw evidence and then run the fail-closed table generator; they
do not trust a preexisting, manually edited summary.  An absent, malformed, or
inconsistent artifact remains not measured before LaTeX compilation.  The
repository-level commands, exact availability, and most recent observed
outcomes are machine-inspected in Table~\ref{tab:artifact-status}.  Command
presence is not evidence that a campaign passed.


\begin{table*}[t]
\centering
\caption{Generated artifact-interface status. ``Present'' means a Make target was found, not that it passed.}
\label{tab:artifact-status}
\footnotesize
\begin{tabular}{llll}
\toprule
Command & Present & Recorded outcome & Source-backed evidence \\
\midrule
\code{make verify} & yes & \notmeasured & no source-backed outcome \\
\code{make formal} & yes & passed (6 models) & formal/results/tlc.log, formal/results/vector\_budget.log, formal/results/sql\_authorization.log, formal/results/multi\_task\_audit.log, formal/results/receipt\_audit.log, formal/results/recovery\_liveness.log \\
\code{make eval-smoke} & yes & \notmeasured & no source-backed outcome \\
\code{make eval-full} & yes & \notmeasured & no source-backed outcome \\
\code{make paper} & yes & \notmeasured & no source-backed outcome \\
\bottomrule
\end{tabular}
\\[1mm]\scriptsize Summary state: source-backed summary loaded.
\end{table*}


\subsection{Claim--Evidence Map}

Table~\ref{tab:claims} distinguishes mechanism existence, test/model evidence,
and claims deliberately withheld.  Paths are relative to the repository root.

\begin{table*}[t]
\centering
\caption{Claim--evidence mapping for this draft.  ``Campaign pending'' means
the source or test exists but no complete raw campaign is asserted here.}
\label{tab:claims}
\scriptsize
\begin{tabular}{L{0.24\textwidth}L{0.33\textwidth}L{0.34\textwidth}}
\toprule
Scoped claim & Implementation / specification evidence & Validation evidence and current restraint \\
\midrule
Manifest identity and digest binding; approval can only narrow & \code{internal/domain/authorization.go}, \code{internal/approval/protocol.go}, OA callback validation & Unit vectors and narrowing/callback tests; cross-component campaign pending. \\
Executing SQL stays within product, column, function/operator, scope, and row envelope & \code{internal/sqlpolicy}, gateway rewrite path, restricted business role & Policy unit corpus and PostgreSQL integration tests; extended adversarial/fuzz campaign pending. \\
$used+reserved\leq limit$ and one in-flight query per task & \code{internal/control/budget.go}, task/budget row locks, TLA+ \code{BudgetSafety}, \code{VectorBudgetSafety}, and \code{SingleInFlight} & Store concurrency tests; finite TLC status is generated separately and is not a general proof. \\
Automatic retry does not execute a request identifier twice & unique $(task\_id,request\_id)$ migration and durable request state & Equal/different-digest replay tests plus \code{AtMostOnceExecution}; a new identifier is explicitly a new operation. \\
Unknown post-boundary outcomes are not released & \code{MarkIndeterminate}, startup recovery, full-reservation charge & Failure-path tests; exhaustive injected crash matrix pending. \\
Catalog/reporting-view drift fails readiness & connector schema attestation at startup and readiness & Connector drift tests; trusted catalog/view policy correctness remains outside scope. \\
Approval/query receipts authenticate service-issued bytes & \code{internal/approval}, \code{internal/queryreceipt}, and gateway terminal-receipt construction & Tamper/key-id vectors and gateway access tests; no claim of human legal non-repudiation or independent DB attestation. \\
Hash chaining makes modification evident relative to a retained trusted prefix & transactional audit sequence, previous/current hashes, receipt binding & Chain verification tests; no WORM or detection of privileged wholesale rewrite is claimed. \\
Performance is acceptable & No claim in this revision & Full raw four-baseline results absent or shown only if generated; therefore no performance conclusion. \\
Human review is faster or less tiring & No claim and no participant data & Only Eq.~\eqref{eq:approval-count} is reported; a human study would require ethics review and real participants. \\
\bottomrule
\end{tabular}
\end{table*}

% Keep the evaluation evidence tables with the evaluation section in both the
% IEEE class and the minimal two-column fallback; otherwise wide floats may
% interrupt the bibliography several pages later.
\clearpage

\section{Limitations}
\label{sec:limitations}

\system's narrowness is part of its security argument and should not be hidden
behind a generic ``agent authorization'' label.

\begin{itemize}
  \item \textbf{No objective semantics.}  The gateway records and displays
  natural language but cannot determine whether a query is necessary, useful,
  or faithful to it.  A human can approve a misleading objective, and allowed
  data can prompt-inject the agent after release.
  \item \textbf{No privacy budget.}  Query count, returned rows, and database
  milliseconds bound resource use, not information leakage.  Repeated allowed
  aggregates may support inference.  There is no differential privacy,
  $k$-anonymity, purpose-use monitoring after release, or output-channel control.
  \item \textbf{Read-only, restricted PostgreSQL only.}  The system supports a
  deny-by-default subset of PostgreSQL~16 \code{SELECT}.  It supports neither
  writes nor arbitrary SQL, stored procedures, a general PG Wire proxy, a
  PostgreSQL extension, payments, or external side effects.
  \item \textbf{Single gateway and per-task serialization.}  There is no
  cross-gateway execution lease or replicated capability store.  Horizontal
  gateway scaling would invalidate the current one-in-flight assumption until a
  shared fencing protocol is added.  Within one gateway, serialization may
  constrain tasks that could otherwise issue parallel reads.
  \item \textbf{Revocation is prospective.}  Revocation prevents a new
  reservation but does not immediately cancel an already in-flight query.  The
  latter is bounded only by its approved statement timeout and grant expiry.
  \item \textbf{Service signatures, not legal non-repudiation.}  An OA
  Ed25519 signature means the trusted OA service recorded a decision attributed
  to an authenticated actor.  It does not prove that the individual personally
  controlled the key, understood the UI, or supplied legally binding intent.
  The gateway signature analogously authenticates a gateway record, not an
  independent execution witness.
  \item \textbf{Unanchored audit.}  Hash chaining detects inconsistencies only
  relative to evidence retained outside the mutable history.  The prototype can
  post signed checkpoint anchors to an external sink, but it does not itself
  provide a WORM replica, transparency log, trusted timestamp, or independent
  retention guarantees for that sink.
  \item \textbf{Trusted policy and infrastructure.}  A malicious/incorrect
  catalog, reporting view, OA, gateway, database administrator, host, or key
  manager defeats the corresponding guarantee.  Schema attestation binds generic
  column shape and normalized PostgreSQL view definitions, but it does not prove
  that a view definition is semantically correct for the intended business
  policy.  Its V1 type check uses generic
  \code{information\_schema.data\_type}; precision, scale, length, domains, and
  array element types are not attested.
  \item \textbf{Prototype identity and key management.}  Static tokens, local
  plaintext HTTP, environment-variable secrets, and demo OA persistence are not
  production controls.  Deployment needs OIDC/OAuth, TLS/mTLS, a durable OA
  outbox, centralized KMS/HSM-backed key custody, result-key material
  destruction, revocation transparency, enterprise WORM or transparency-log
  retention for audit anchors,
  backup/restore, organizational retention workflows, rate limits, and
  operational monitoring.
  \item \textbf{Incomplete empirical evidence.}  A finite TLC run does not
  prove unbounded behavior, fuzzing cannot prove parser-policy completeness, and
  this draft makes no performance or human-factors conclusion without the raw
  campaigns described in Sec.~\ref{sec:evaluation}.
\end{itemize}

\section{Conclusion}
\label{sec:conclusion}

\system demonstrates a concrete way to authorize a bounded family of future
database reads without handing an autonomous agent a reusable database
credential.  A human-visible manifest, narrowing-only grant, deterministic SQL
enforcement, transactional reservation ledger, explicit indeterminate state,
schema attestation, and scoped receipts together address authorization drift,
retry, concurrency, and crash ambiguity in the fixed single-gateway PostgreSQL
deployment.  The mechanism is best understood as task-bound authorization and
continuous usage control, not as a new authentication framework or a semantic
alignment oracle.  Its strongest conclusions remain correspondingly modest:
the design makes a precise enforcement boundary implementable and auditable;
performance, broad SQL coverage, distributed deployment, privacy protection,
and human-review benefits require evidence and mechanisms not claimed here.

\IfFileExists{IEEEtran.bst}{\bibliographystyle{IEEEtran}}{\bibliographystyle{plain}}
\bibliography{references}

\end{document}
